\chapter{Design and Trade-offs}

A language is a set of decisions, and decisions have costs as well as
benefits. This chapter argues for wick's decisions honestly: it puts the wick
and Lua builds of the same real game side by side, it says plainly where Lua
still reads better, and it explains the roadmap philosophy that makes v0.1's
omissions deliberate rather than embarrassing.

\section{The same game, twice}

Kora Night is the proof. It ships in both languages, as two deliberately
separate projects --- `games/showcase` in Lua and `games/showcase_wick` in
wick --- kept side by side so the two languages stay comparable on the same
real game. Read the pair like a Rosetta stone. Three excerpts carry the
argument.

\subsection{The hi-score load}

This is the bug from Chapter~4, in its natural habitat. The Lua build shipped
with a frame-1 crash:

\begin{lstlisting}[language=luatwin]
-- Lua: a 0-byte save makes load_save return "" (truthy!), the
-- fallback never fires, and tonumber("") = nil poisons everything.
local best = tonumber(lt.load_save("koranight_best") or "0")
\end{lstlisting}

\noindent
The wick build cannot express that bug. The line below is the only version
that compiles, because both ways the value can be absent --- no file, empty
file --- are forced to be handled:

\begin{lstlisting}[language=wick]
// wick: the ONLY version the compiler accepts.
let best = num(lt.load_save("koranight_wick_best") ?? "") ?? 0
\end{lstlisting}

\noindent
Here wick wins outright. The safe code is shorter than the buggy Lua code, and
the unsafe code does not compile.

\subsection{Records}

Now the case wick loses. A lamp has a position and a lit flag. In Lua it is a
table with named fields, and a list of lamps is a list of tables:

\begin{lstlisting}[language=luatwin]
-- Lua: a list of records, read naturally
local LAMPS = {
  { x = 3.4, z = 3.4, lit = true },
  -- ...
}
for i, l in ipairs(LAMPS) do  -- l.lit, l.x, l.z
end
\end{lstlisting}

\noindent
wick v0.1 has no structs, so the same data becomes typed parallel lists, with
an index standing in for a record reference:

\begin{lstlisting}[language=wick]
// wick v0.1: typed parallel lists (structs are on the roadmap)
let lamp_x: list<num> = []
let lamp_z: list<num> = []
let lamp_lit: list<bool> = []
for i in 0..4 {
  // lamp_lit[i], lamp_x[i], lamp_z[i]
}
\end{lstlisting}

\noindent
The honest score: Lua reads better here today. A lamp is one thing, and Lua
lets you write it as one thing; wick scatters it across three lists you must
keep in step. wick's version is fully typed, which Lua's is not --- `l.lit` is
whatever it happens to be, where `lamp_lit[i]` is a checked `bool` --- but that
is small consolation for the loss of readability. Records are the top of
wick's v0.2 wishlist for exactly this reason.

\subsection{Multiple returns}

A smaller loss. Lua returns the nearest lamp and its distance together:

\begin{lstlisting}[language=luatwin]
local i, dist2 = nearest_unlit()
\end{lstlisting}

\noindent
wick has no multiple returns, so the two queries are split, and ``no lamp''
becomes a sentinel index rather than an absent second value:

\begin{lstlisting}[language=wick]
let i = nearest_unlit()      // returns the index, or -1
let d = dist2_to(i)          // second query, explicit
\end{lstlisting}

\noindent
This is more verbose, but it is not worse in kind --- the split is mechanical,
and the sentinel is a well-understood idiom. Multiple returns are on the
roadmap; their absence is an inconvenience, not a hazard.

\section{What you stop writing}

Set against those losses is a real reduction in the defensive code a dynamic
language demands. Porting Kora Night to wick, the following simply left the
source:

\begin{itemize}
\item The `if x ~= nil` chains around every save and parse. The compiler
tracks nullability, so you write `??` once at the boundary and the value is a
plain `T` thereafter.
\item The `math.` prefixes. `sin`, `floor`, `rand`, and their kin are
built-ins with no namespace.
\item The `string.format` calls that existed only to print an integer without
a trailing decimal. `str(n)` prints integers cleanly.
\item The argument-order paranoia on engine calls. Arity and types are checked
at compile time with the call name and the argument number in the message, so
a mistake is caught the moment it is written, not discovered by a wrong-looking
frame.
\end{itemize}

\section{What stays the same}

It is worth being clear about what the language change did \emph{not} disturb,
because it bounds the cost of adopting wick. The engine API is call-for-call
identical: `lt.draw(...)` is the same line in both builds. The frame contract
--- `update(dt)` and `draw()` --- is the same. Hot reload and the error screen
behave the same. And the two games render the same frames, which is not a
figure of speech: the deterministic screenshot tests of Chapter~6 hold the
wick and Lua builds to the same pixels. wick changed the language, not the
engine and not the workflow.

\section{The v0.1 limits}

wick is deliberately small --- Lua-1.0 small. Every omission below is a known
gap with a workaround today and a place on the roadmap, not a surprise waiting
to be discovered. The bar for v0.1 was fixed and narrow: run real lantern
games, kill the Lua bug classes, and stay around two thousand lines.

\begin{center}
\begin{tabular}{@{}p{0.42\linewidth}p{0.50\linewidth}@{}}
\toprule
\textbf{Not in v0.1} & \textbf{Workaround today} \\
\midrule
Structs / records & typed parallel lists (see the port above) \\
Closures, nested and first-class functions & top-level \texttt{fn} plus module globals \\
Multiple return values & two functions, or out-of-band state \\
Nested containers (\texttt{list<list<num>>}) & index arithmetic on a flat list \\
Modules / imports (one \texttt{main.wick}) & it is a game script, not an app platform --- yet \\
\texttt{for x in list} iteration & \texttt{for i in 0..len(xs)} \\
String indexing / slicing / methods & \texttt{len}, \texttt{+}, \texttt{str()} cover the HUD cases \\
Exhaustive return-path checking & a typed function that falls off the end returns \texttt{nil} --- don't \\
Narrowing for globals / \texttt{and}-chains & \texttt{??}, or copy to a local first \\
Hex / exponent number literals & decimal \\
\bottomrule
\end{tabular}
\end{center}

\noindent
Two entries deserve to be pulled out, because they are not limits at all but
choices, and they are choices the roadmap will not reverse:

\paragraph{No JIT.} This is iOS-legal by construction and gives deterministic
performance with no warm-up. It is a feature, permanent.

\paragraph{No time-seeded randomness.} Determinism is the default so that
replays and CI screenshots work; if you want variety between runs you seed
`srand()` yourself, with a number you can log. Also permanent.

\section{Roadmap philosophy}

The rule that governs what happens next deserves stating as law, because
it is the core thing this language follows: \emph{wick contains only what
we truly need --- nothing for the sake of having it.} A feature is
admitted when real game code demonstrably hurt without it, and not
before. The paired Kora Night builds are the standing evidence mechanism:
missing-feature pain must show up there, in code anyone can read, before
the language grows. Optionals are in this book because a shipped crash
demanded them; closures are not, because no shipped program has yet paid
a price for their absence. Feature parity with other languages is
explicitly not a reason to add anything.


The organising idea is ``Lua 1.0 small.'' Lua did not begin with metatables,
coroutines, and an ecosystem; it began as a small, clean, embeddable core and
grew as real use justified each addition. wick is taking the same path
deliberately. Every omission in the table above is addable behind the same
typed front end that already exists --- records, closures, and multiple returns
do not require rethinking the compiler, only extending it --- and each will be
added when a real game needs it enough to justify the size, not before.

This is why the limits are stated so plainly. A language that hid its gaps
would be inviting you to discover them at the worst moment. wick lists them,
gives each a workaround, and puts each on a roadmap, because the whole thesis
of the language is that a program's failure modes should be visible up front
rather than discovered in production --- and a language owes its users the same
honesty it enforces in their code.
